\documentclass[../../../include/open-logic-section]{subfiles}

\begin{document}

\olfileid[tr]{sth}{replacement}{ref}
\olsection{Yerine Koyma ve Yansıma}

Yerine Koyma'yı \stagesinex{} aracılığıyla gerekçelendirmeye yönelik son
girişimimiz, derin ve güzel bir sonuçla başlar:\footnote{Hatırlatma: Açıkça
aksi belirtilmedikçe bütün formüller parametre içerebilir.}
%In this, I will use overlining, such as $x_1, \ldots, x_n$, to
%abbreviate ``$x_1, \ldots, x_n$'': 
\begin{thm}[Yansıma Şeması]\ollabel{reflectionschema}
Her $\phi$ formülü için:
\[
\forall \alpha \exists \beta > \alpha (\forall x_1 \ldots, x_n \in
V_\beta)(\phi(x_1, \ldots, x_n) \liff \phi^{V_\beta}(x_1, \ldots, x_n))
\]
\end{thm}
\noindent 
\olref[sth][replacement][strength]{formularelativization} içinde olduğu gibi
$\phi^{V_\beta}$, $\phi$ içindeki her niceleyicinin $V_\beta$ kümesiyle
sınırlandırılması sonucudur. Dolayısıyla Yansıma sezgisel olarak şunu söyler:
$\phi$ bütün hiyerarşide doğruysa $\phi$ hiyerarşinin keyfî ölçüde çok
\emph{başlangıç kesiminde} de doğrudur.

\citet{Montague1961} ile \citet{Levy1960}, Yerine Koyma ve Yansıma'nın (uygun
formülasyonlarının) $\Z$ üzerinden denk olduğunu, dolayısıyla ikisinden
birinin eklenmesiyle $\ZF$'nin elde edildiğini gösterdi. (Bu sonuçları
\olref[sth][replacement][refproofs]{sec} içinde ispatlıyoruz.) Bu denklik göz
önünde bulundurulunca Yansıma ve Yerine Koyma'yı \stagesinex{} aracılığıyla
şöyle gerekçelendirmeyi umabiliriz: \stagesinex{} verildiğinde hiyerarşi çok,
çok yüksek; hatta onun hakkında söyleyebileceğimiz hiçbir şey yüksekliğini
sınırlandırmaya yetmeyecek kadar yüksek olmalıdır. Bunu, herhangi bir
$\phi$ tümcesi bütün hiyerarşide doğruysa hiyerarşinin keyfî ölçüde çok
başlangıç kesiminde de doğru olduğu düşüncesi olarak anlayabiliriz. Bu da tam
olarak Yansıma'dır.

Bu da Yerine Koyma için içsel bir gerekçelendirme sağlamaya yönelik gerçekten
umut verici bir girişim gibi görünür. Fakat burada söylenebileceklerin miktarı
çok fazladır. Başarılı olup olmadığına artık kendiniz karar vermelisiniz.\footnote{Yine de
\citet[95--100]{Incurvati2020} okuyarak devam etmek isteyebilirsiniz.}

\end{document}
